\begin{abstract}
Discrete visual abstraction benchmarks like the Abstraction and Reasoning Corpus (ARC) require autonomous agents to infer generative rules from few demonstration pairs without leaking ungrounded search priors into downstream execution. In this work, we introduce \texttt{NeuroARC.jl}, a high performance Julia framework that operationalizes Peter Tse's cognitive theory of criterial causation, Wood-Bruner-Ross pedagogical scaffolding, and typed hyphal grafting within an auditable execution pipeline. Rather than treating candidate selection as unconstrained neural generation or unstratified brute-force program search, \texttt{NeuroARC.jl} formalizes rule composition as a tri-stage temporal process ($t_1, t_2, t_3$) governed by criterial logic gates (Threshold, Bandpass, Subtype, Conjunction, Disjunction, and K-of-N). Candidate transformations are composed via static subtyping preflight verification ($\mathrm{Out}(A) <: \mathrm{In}(B)$) paired with runtime boundary enforcement, preventing numeric sentinel leakage by emitting algebraic verdicts (\texttt{Accepted}, \texttt{Rejected}, and \texttt{Incomplete}). We isolate a zero-dependency core package while providing explicit extension boundaries for Apple Silicon GPU acceleration via \texttt{Metal.jl}, cross-architecture parallel computing via \texttt{JACC.jl}, and probabilistic calibration via \texttt{Turing.jl}. This package structure is accompanied by an automated continuous integration pipeline that compiles this manuscript and validates 325 unit tests under clean runners.
\end{abstract}
